% --- LaTeX Homework Template - S. Venkatraman ---

% --- Set document class and font size ---

\documentclass[letterpaper, 11pt]{article}

% --- Package imports ---

\usepackage{
  amsmath, amsthm, amssymb, mathtools, dsfont,	  % Math typesetting
  graphicx, wrapfig, subfig, float,                  % Figures and graphics formatting
  listings, color, inconsolata, pythonhighlight,     % Code formatting
  fancyhdr, sectsty, hyperref, enumerate, enumitem } % Headers/footers, section fonts, links, lists

% --- Page layout settings ---
\usepackage{bm}
% Set page margins
\usepackage[left=1.35in, right=1.35in, bottom=1in, top=1.1in, headsep=0.2in]{geometry}

% Anchor footnotes to the bottom of the page
\usepackage[bottom]{footmisc}

% Set line spacing
\renewcommand{\baselinestretch}{1.2}

% Set spacing between paragraphs
\setlength{\parskip}{1.5mm}

% Allow multi-line equations to break onto the next page
\allowdisplaybreaks

% Enumerated lists: make numbers flush left, with parentheses around them
\setlist[enumerate]{wide=0pt, leftmargin=21pt, labelwidth=0pt, align=left}
\setenumerate[1]{label={(\arabic*)}}

% --- Page formatting settings ---

% Set link colors for labeled items (blue) and citations (red)
\hypersetup{colorlinks=true, linkcolor=blue, citecolor=red}

% Make reference section title font smaller
\renewcommand{\refname}{\large\bf{References}}

% --- Settings for printing computer code ---

% Define colors for green text (comments), grey text (line numbers),
% and green frame around code
\definecolor{greenText}{rgb}{0.5, 0.7, 0.5}
\definecolor{greyText}{rgb}{0.5, 0.5, 0.5}
\definecolor{codeFrame}{rgb}{0.5, 0.7, 0.5}

% Define code settings
\lstdefinestyle{code} {
  frame=single, rulecolor=\color{codeFrame},            % Include a green frame around the code
  numbers=left,                                         % Include line numbers
  numbersep=8pt,                                        % Add space between line numbers and frame
  numberstyle=\tiny\color{greyText},                    % Line number font size (tiny) and color (grey)
  commentstyle=\color{greenText},                       % Put comments in green text
  basicstyle=\linespread{1.1}\ttfamily\footnotesize,    % Set code line spacing
  keywordstyle=\ttfamily\footnotesize,                  % No special formatting for keywords
  showstringspaces=false,                               % No marks for spaces
  xleftmargin=1.95em,                                   % Align code frame with main text
  framexleftmargin=1.6em,                               % Extend frame left margin to include line numbers
  breaklines=true,                                      % Wrap long lines of code
  postbreak=\mbox{\textcolor{greenText}{$\hookrightarrow$}\space} % Mark wrapped lines with an arrow
}

% Set all code listings to be styled with the above settings
\lstset{style=code}

% --- Math/Statistics commands ---

% Add a reference number to a single line of a multi-line equation
% Usage: "\numberthis\label{labelNameHere}" in an align or gather environment
\newcommand\numberthis{\addtocounter{equation}{1}\tag{\theequation}}

% Shortcut for bold text in math mode, e.g. $\b{X}$
\let\b\mathbf

% Shortcut for bold Greek letters, e.g. $\bg{\beta}$
\let\bg\boldsymbol

% Shortcut for calligraphic script, e.g. %\mc{M}$
\let\mc\mathcal

% \mathscr{(letter here)} is sometimes used to denote vector spaces
\usepackage[mathscr]{euscript}

% Convergence: right arrow with optional text on top
% E.g. $\converge[w]$ for weak convergence
\newcommand{\converge}[1][]{\xrightarrow{#1}}

% Normal distribution: arguments are the mean and variance
% E.g. $\normal{\mu}{\sigma}$
\newcommand{\normal}[2]{\mathcal{N}\left(#1,#2\right)}

% Uniform distribution: arguments are the left and right endpoints
% E.g. $\unif{0}{1}$
\newcommand{\unif}[2]{\text{Uniform}(#1,#2)}

% Independent and identically distributed random variables
% E.g. $ X_1,...,X_n \iid \normal{0}{1}$
\newcommand{\iid}{\stackrel{\smash{\text{iid}}}{\sim}}

% Equality: equals sign with optional text on top
% E.g. $X \equals[d] Y$ for equality in distribution
\newcommand{\equals}[1][]{\stackrel{\smash{#1}}{=}}

% Math mode symbols for common sets and spaces. Example usage: $\R$
\newcommand{\R}{\mathbb{R}}   % Real numbers
\newcommand{\C}{\mathbb{C}}   % Complex numbers
\newcommand{\Q}{\mathbb{Q}}   % Rational numbers
\newcommand{\Z}{\mathbb{Z}}   % Integers
\newcommand{\N}{\mathbb{N}}   % Natural numbers
\newcommand{\F}{\mathcal{F}}  % Calligraphic F for a sigma algebra
\newcommand{\El}{\mathcal{L}} % Calligraphic L, e.g. for L^p spaces

% Math mode symbols for probability
\newcommand{\pr}{\mathbb{P}}    % Probability measure
\newcommand{\E}{\mathbb{E}}     % Expectation, e.g. $\E(X)$
\newcommand{\var}{\text{Var}}   % Variance, e.g. $\var(X)$
\newcommand{\cov}{\text{Cov}}   % Covariance, e.g. $\cov(X,Y)$
\newcommand{\corr}{\text{Corr}} % Correlation, e.g. $\corr(X,Y)$
\newcommand{\B}{\mathcal{B}}    % Borel sigma-algebra

% Other miscellaneous symbols
\newcommand{\tth}{\text{th}}	% Non-italicized 'th', e.g. $n^\tth$
\newcommand{\Oh}{\mathcal{O}}	% Big-O notation, e.g. $\O(n)$
\newcommand{\1}{\mathds{1}}	% Indicator function, e.g. $\1_A$

% Additional commands for math mode
\DeclareMathOperator*{\argmax}{argmax}    % Argmax, e.g. $\argmax_{x\in[0,1]} f(x)$
\DeclareMathOperator*{\argmin}{argmin}    % Argmin, e.g. $\argmin_{x\in[0,1]} f(x)$
\DeclareMathOperator*{\spann}{Span}       % Span, e.g. $\spann\{X_1,...,X_n\}$
\DeclareMathOperator*{\bias}{Bias}        % Bias, e.g. $\bias(\hat\theta)$
\DeclareMathOperator*{\ran}{ran}          % Range of an operator, e.g. $\ran(T) 
\DeclareMathOperator*{\dv}{d\!}           % Non-italicized 'with respect to', e.g. $\int f(x) \dv x$
\DeclareMathOperator*{\diag}{diag}        % Diagonal of a matrix, e.g. $\diag(M)$
\DeclareMathOperator*{\trace}{trace}      % Trace of a matrix, e.g. $\trace(M)$

% Numbered theorem, lemma, etc. settings - e.g., a definition, lemma, and theorem appearing in that 
% order in Section 2 will be numbered Definition 2.1, Lemma 2.2, Theorem 2.3. 
% Example usage: \begin{theorem}[Name of theorem] Theorem statement \end{theorem}
\theoremstyle{definition}
\newtheorem{theorem}{Theorem}[section]
\newtheorem{proposition}[theorem]{Proposition}
\newtheorem{lemma}[theorem]{Lemma}
\newtheorem{corollary}[theorem]{Corollary}
\newtheorem{definition}[theorem]{Definition}
\newtheorem{example}[theorem]{Example}
\newtheorem{remark}[theorem]{Remark}

% Un-numbered theorem, lemma, etc. settings
% Example usage: \begin{lemma*}[Name of lemma] Lemma statement \end{lemma*}
\newtheorem*{theorem*}{Theorem}
\newtheorem*{proposition*}{Proposition}
\newtheorem*{lemma*}{Lemma}
\newtheorem*{corollary*}{Corollary}
\newtheorem*{definition*}{Definition}
\newtheorem*{example*}{Example}
\newtheorem*{remark*}{Remark}
\newtheorem*{claim}{Claim}

\newcommand{\Pp}{\mathbb{P}}
\newcommand{\Dkl}{\mathrm{D}}

% --- Left/right header text (to appear on every page) ---

% Include a line underneath the header, no footer line
\pagestyle{fancy}
\renewcommand{\footrulewidth}{0pt}
\renewcommand{\headrulewidth}{0.4pt}

% Left header text: course name/assignment number
\lhead{Foundations of Machine Learning-- Homework 1}

% Right header text: your name
\rhead{Yixia Yu (yy5091@nyu.edu)}

% --- Document starts here ---

\begin{document}
\section*{A. VC-dimension}

\begin{enumerate}
\item Let $\mathcal{G}$ be a finite-dimensional vector space of real-valued functions on $\mathbb{R}^n$, with $\dim(\mathcal{G}) = k < +\infty$. Define the class of sets:
\[
\mathcal{H} = \left\{\left\{x \in \mathbb{R}^n: g(x) \geq 0\right\}: g \in \mathcal{G}\right\}.
\]

Show that the VC-dimension $d$ of $\mathcal{H}$ is finite and satisfies $d \leq k$. [\textit{Hint:} Consider an arbitrary collection of $m = k + 1$ points and the linear map $T: \mathcal{G} \to \mathbb{R}^m$ defined by
\[
T(g) = (g(y_1), \ldots, g(y_m)).
\]

\item Let $\mathcal{F}$ be a finite-dimensional vector space of real functions on $\mathbb{R}^n$, with $\dim(\mathcal{F}) = r < +\infty$. Consider the set of hypotheses
\[
H = \left\{x \in \mathbb{R}^n: \sum_{j=1}^k a_j f_j(x) \geq 0\right\}: f_1, \ldots, f_k \in \mathcal{F}, \ a_1, \ldots, a_k \in \mathbb{R}.
\]

Show that the VC dimension $d$ of $H$ is finite and satisfies $d \leq kr$. [\textit{Hint:} Map each $x_i$ to a vector in $\mathbb{R}^{kr}$ representing all evaluations of the $k$ functions, then use the standard result that the VC dimension of linear threshold functions in $\mathbb{R}^m$ is $m$.]

\item Let $\mathcal{F}$ be a finite-dimensional vector space of real functions on $\mathbb{R}^n$ with $\dim(\mathcal{F}) = r < +\infty$. Consider the class of sets defined by intersections of thresholded functions from $\mathcal{F}$:
\[
H = \left\{\bigcap_{j=1}^k \left\{x \in \mathbb{R}^n: f_j(x) \geq 0\right\}: f_1, \ldots, f_k \in \mathcal{F}\right\}.
\]

That is, each hypothesis selects the points where all $k$ functions are non-negative simultaneously.

Show that the VC dimension $d$ of $H$ is finite, and provide a bound on $d$ in terms of $k$ and $r$. [\textit{Hint:} You can prove that the VC-dimension of intersections of $k$ classes of finite VC dimension $r$ can be bounded by $d \leq 2kr \log_2(2k)$. Consider using Sauer's lemma iteratively for the intersection and mapping the problem to a high-dimensional linear space if necessary.]
\end{enumerate}

\section*{Solutions}

(1)

\textbf{Proof:}

We prove that the VC-dimension $d$ of $\mathcal{H}$ satisfies $d \leq k$ by showing that no set of $m = k + 1$ points can be shattered by $\mathcal{H}$.

Let $\{y_1, y_2, \ldots, y_m\}$ be an arbitrary collection of $m = k + 1$ points in $\mathbb{R}^n$.

Define the linear map $T: \mathcal{G} \to \mathbb{R}^m$ by
\[
T(g) = (g(y_1), g(y_2), \ldots, g(y_m)).
\]

Since $\mathcal{G}$ is a vector space of dimension $k$ and $T$ is a linear map, by the rank-nullity theorem, we have
\[
\dim(\text{Im}(T)) \leq \dim(\mathcal{G}) = k < m = k + 1.
\]

Therefore, $\text{Im}(T)$ is a proper subspace of $\mathbb{R}^m$ with dimension at most $k < m$.

This means there exists a non-zero vector $v = (v_1, v_2, \ldots, v_m) \in \mathbb{R}^m$ that is orthogonal to $\text{Im}(T)$. That is, for all $g \in \mathcal{G}$,
\[
\langle T(g), v \rangle = \sum_{i=1}^m v_i g(y_i) = 0.
\]

Since $v \neq 0$, we can partition the indices $\{1, 2, \ldots, m\}$ into three disjoint sets:
\begin{align*}
I^+ &= \{i : v_i > 0\}, \\
I^- &= \{i : v_i < 0\}, \\
I^0 &= \{i : v_i = 0\}.
\end{align*}

Note that $I^+ \cup I^- \neq \emptyset$ since $v \neq 0$.

We claim that if $I^+ \neq \emptyset$ and $I^- \neq \emptyset$, then the dichotomy corresponding to the subset $I^+$ (i.e., classifying points in $I^+$ as positive and points in $I^- \cup I^0$ as negative) cannot be realized by any hypothesis in $\mathcal{H}$.

Suppose for contradiction that there exists $g \in \mathcal{G}$ such that 
\begin{align*}
g(y_i) &\geq 0 \text{ for all } i \in I^+, \\
g(y_i) &< 0 \text{ for all } i \in I^-.
\end{align*}

Then:
\[
\sum_{i=1}^m v_i g(y_i) = \sum_{i \in I^+} v_i g(y_i) + \sum_{i \in I^-} v_i g(y_i) + \sum_{i \in I^0} v_i g(y_i).
\]

Since $v_i = 0$ for $i \in I^0$, the third sum vanishes.

For the first sum: since $v_i > 0$ and $g(y_i) \geq 0$ for $i \in I^+$, we have $\sum_{i \in I^+} v_i g(y_i) \geq 0$.

For the second sum: since $v_i < 0$ and $g(y_i) < 0$ for $i \in I^-$, we have $v_i g(y_i) > 0$, so $\sum_{i \in I^-} v_i g(y_i) > 0$ (assuming $I^- \neq \emptyset$).

Therefore, $\sum_{i=1}^m v_i g(y_i) > 0$, which contradicts the orthogonality condition.

If $I^- = \emptyset$ but $I^+ \neq \emptyset$, a similar argument shows that no $g \in \mathcal{G}$ can satisfy $g(y_i) < 0$ for all $i \in I^+$, which means the empty subset cannot be realized (contradiction to shattering).

Thus, the collection $\{y_1, \ldots, y_m\}$ cannot be shattered by $\mathcal{H}$.

Since no set of $k + 1$ points can be shattered, we conclude that the VC-dimension $d$ of $\mathcal{H}$ satisfies $d \leq k$. \qed


(2)

\textbf{Proof:}

We show that $H$ can be viewed as a class of linear threshold functions (homogeneous half-spaces) in a finite-dimensional space, from which the bound follows.

Let $\{f_1^*, \ldots, f_r^*\}$ be a basis for $\mathcal{F}$. Any function $f \in \mathcal{F}$ can be written as $f = \sum_{\ell=1}^r \beta_\ell f_\ell^*$ for some coefficients $\beta_1, \ldots, \beta_r \in \mathbb{R}$.

Consider the feature map $\Phi: \mathbb{R}^n \to \mathbb{R}^{kr}$ defined by
\[
\Phi(x) = (f_1^*(x), \ldots, f_r^*(x), f_1^*(x), \ldots, f_r^*(x), \ldots, f_1^*(x), \ldots, f_r^*(x)),
\]
where the tuple $(f_1^*(x), \ldots, f_r^*(x))$ is repeated $k$ times, arranged in $k$ blocks of size $r$.

For any hypothesis in $H$ defined by functions $f_1, \ldots, f_k \in \mathcal{F}$ and coefficients $a_1, \ldots, a_k \in \mathbb{R}$, where $f_j = \sum_{\ell=1}^r \beta_{j\ell} f_\ell^*$, the decision rule is:
\[
\sum_{j=1}^k a_j f_j(x) = \sum_{j=1}^k a_j \sum_{\ell=1}^r \beta_{j\ell} f_\ell^*(x) \geq 0.
\]

This can be rewritten as:
\[
\langle w, \Phi(x) \rangle \geq 0,
\]
where the weight vector $w \in \mathbb{R}^{kr}$ is given by
\[
w = (a_1\beta_{1,1}, \ldots, a_1\beta_{1,r}, a_2\beta_{2,1}, \ldots, a_2\beta_{2,r}, \ldots, a_k\beta_{k,1}, \ldots, a_k\beta_{k,r}).
\]

Therefore, $H$ is equivalent to the class of homogeneous linear threshold functions in $\mathbb{R}^{kr}$:
\[
H \cong \{\{z \in \mathbb{R}^{kr} : w \cdot z \geq 0\} : w \in \mathbb{R}^{kr}\}.
\]

It is a well-known result that the VC-dimension of homogeneous half-spaces (linear threshold functions through the origin) in $\mathbb{R}^m$ is exactly $m$.

Therefore, the VC-dimension $d$ of $H$ satisfies $d \leq kr$.


(3)

\textbf{Proof:}

We prove that the VC-dimension of $H$ is finite and provide a bound in terms of $k$ and $r$.

First, note that by Problem 1, the class
\[
\mathcal{H}_0 = \{\{x \in \mathbb{R}^n : f(x) \geq 0\} : f \in \mathcal{F}\}
\]
has VC-dimension at most $r$.

The class $H$ consists of intersections of $k$ sets from $\mathcal{H}_0$:
\[
H = \left\{\bigcap_{j=1}^k S_j : S_1, \ldots, S_k \in \mathcal{H}_0\right\}.
\]

We use the following general result on intersections:

\textbf{Lemma (VC-dimension of intersections):} Let $\mathcal{C}_1, \ldots, \mathcal{C}_k$ be hypothesis classes, each with VC-dimension at most $d$. Then the class of intersections
\[
\mathcal{C} = \left\{\bigcap_{i=1}^k C_i : C_i \in \mathcal{C}_i\right\}
\]
has VC-dimension at most $O(kd \log k)$.

More precisely, it can be shown that:
\[
\text{VC-dim}(\mathcal{C}) \leq 2kd \log_2(3k).
\]

\textbf{Sketch of proof of the lemma:}

Let $S = \{x_1, \ldots, x_m\}$ be a set of $m$ points. The growth function $\Pi_{\mathcal{C}}(m)$ counts the number of distinct ways the class $\mathcal{C}$ can label the points in $S$.

For the intersection $\bigcap_{i=1}^k C_i$ to produce a specific labeling, each $C_i$ must produce a labeling compatible with the intersection. By Sauer's lemma, each class $\mathcal{C}_i$ with VC-dimension $d$ can produce at most $\sum_{j=0}^d \binom{m}{j} \leq \left(\frac{em}{d}\right)^d$ distinct labelings on $m$ points.

The key insight is that for the intersection to shatter a set, we need enough "degrees of freedom" from all $k$ classes combined. Using a double counting argument and the constraint that the intersection must be consistent, one can show that if $m > 2kd\log_2(3k)$, then $\Pi_{\mathcal{C}}(m) < 2^m$, which means $S$ cannot be shattered.

\textbf{Application to our problem:}

In our case, we have $k$ copies of the same class $\mathcal{H}_0$, each with VC-dimension at most $r$. Applying the lemma with $d = r$, we obtain:
\[
\text{VC-dim}(H) \leq 2kr \log_2(3k).
\]

Therefore, the VC-dimension $d$ of $H$ is finite and satisfies
\[
d \leq 2kr \log_2(3k) = O(kr \log k).
\]

\section*{B. Rademacher complexity}

\begin{enumerate}
\item Let $G$ be a family of functions mapping from $Z$ to $[0,1]$. The general Rademacher complexity bound presented in class was based on the analysis of the function $\Phi$ defined by $\Phi(S) = \sup_{g \in G} \mathbb{E}[g] - \widehat{\mathbb{E}}_S[g]$ for any training sample $S = (z_1, \ldots, z_m)$ of size $m$, with $\widehat{\mathbb{E}}_S[g] = \frac{1}{m}\sum_{i=1}^m g(z_i)$. Instead, apply McDiarmid's inequality to $\Psi$ defined by $\Psi(S) = \sup_{g \in G} \mathbb{E}[g] - \widehat{\mathbb{E}}_S[g] - 2\widehat{\mathfrak{R}}_S(G)$ and try to obtain a slightly better generalization bound than the one obtained in class in terms of the empirical Rademacher complexity.
\end{enumerate}

Let $S = (z_1, \ldots, z_m)$ and $S^{(i)} = (z_1, \ldots, z_{i-1}, z_i', z_{i+1}, \ldots, z_m)$ differ only in the $i$-th coordinate.

We need to bound:
$$|\Psi(S) - \Psi(S^{(i)})|$$

\textbf{Analysis:}

$$\Psi(S) = \sup_{g \in G} \left[\mathbb{E}[g] - \widehat{\mathbb{E}}_S[g] - 2\widehat{\mathfrak{R}}_S(G)\right]$$

where $\widehat{\mathbb{E}}_S[g] = \frac{1}{m}\sum_{j=1}^m g(z_j)$ and 
$$\widehat{\mathfrak{R}}_S(G) = \mathbb{E}_{\sigma}\left[\sup_{g \in G} \frac{1}{m}\sum_{j=1}^m \sigma_j g(z_j)\right].$$

\textbf{Change in empirical mean:}
$$\widehat{\mathbb{E}}_S[g] - \widehat{\mathbb{E}}_{S^{(i)}}[g] = \frac{1}{m}[g(z_i) - g(z_i')]$$

Since $g \in [0,1]$, we have $|g(z_i) - g(z_i')| \leq 1$.

\textbf{Change in empirical Rademacher complexity:}
\begin{align*}
&\left|\widehat{\mathfrak{R}}_S(G) - \widehat{\mathfrak{R}}_{S^{(i)}}(G)\right| \\
&= \left|\mathbb{E}_{\sigma}\left[\sup_{g \in G} \frac{1}{m}\sum_{j=1}^m \sigma_j g(z_j)\right] - \mathbb{E}_{\sigma}\left[\sup_{g \in G} \frac{1}{m}\sum_{j=1}^m \sigma_j g(z_j^{(i)})\right]\right|
\end{align*}

where $z_j^{(i)} = z_j$ for $j \neq i$ and $z_i^{(i)} = z_i'$.

By the triangle inequality and properties of supremum:
\begin{align*}
\left|\widehat{\mathfrak{R}}_S(G) - \widehat{\mathfrak{R}}_{S^{(i)}}(G)\right| 
&\leq \mathbb{E}_{\sigma}\left[\sup_{g \in G} \frac{1}{m}|\sigma_i||g(z_i) - g(z_i')|\right] \\
&\leq \frac{1}{m}\mathbb{E}_{\sigma}[|\sigma_i|] \cdot 1 = \frac{1}{m}.
\end{align*}

\textbf{Combining both changes:}

For any $g \in G$:
$$\left|\left(\mathbb{E}[g] - \widehat{\mathbb{E}}_S[g] - 2\widehat{\mathfrak{R}}_S(G)\right) - \left(\mathbb{E}[g] - \widehat{\mathbb{E}}_{S^{(i)}}[g] - 2\widehat{\mathfrak{R}}_{S^{(i)}}(G)\right)\right|$$
$$= \left|\widehat{\mathbb{E}}_{S^{(i)}}[g] - \widehat{\mathbb{E}}_S[g] + 2(\widehat{\mathfrak{R}}_{S^{(i)}}(G) - \widehat{\mathfrak{R}}_S(G))\right|$$
$$\leq \frac{1}{m} + 2 \cdot \frac{1}{m} = \frac{3}{m}.$$

Therefore:
$$|\Psi(S) - \Psi(S^{(i)})| \leq \frac{3}{m}.$$

So $c_i = \frac{3}{m}$ for all $i \in [m]$.
\\
Then, we need to prove that
$$\mathbb{E}[\Psi(S)] \leq 0$$

\begin{proof}
By symmetrization (using ghost samples):
\begin{align*}
\mathbb{E}_S[\Psi(S)] &= \mathbb{E}_S\left[\sup_{g \in G} \mathbb{E}[g] - \widehat{\mathbb{E}}_S[g] - 2\widehat{\mathfrak{R}}_S(G)\right] \\
&= \mathbb{E}_S\left[\sup_{g \in G} \mathbb{E}_{S'}[\widehat{\mathbb{E}}_{S'}[g]] - \widehat{\mathbb{E}}_S[g] - 2\widehat{\mathfrak{R}}_S(G)\right] \\
&\leq \mathbb{E}_{S,S'}\left[\sup_{g \in G} (\widehat{\mathbb{E}}_{S'}[g] - \widehat{\mathbb{E}}_S[g])\right] - 2\mathbb{E}_S[\widehat{\mathfrak{R}}_S(G)]
\end{align*}

By the symmetrization lemma:
$$\mathbb{E}_{S,S'}\left[\sup_{g \in G} (\widehat{\mathbb{E}}_{S'}[g] - \widehat{\mathbb{E}}_S[g])\right] \leq 2\mathbb{E}_{S,\sigma}\left[\sup_{g \in G} \frac{1}{m}\sum_{i=1}^m \sigma_i g(z_i)\right] = 2\mathbb{E}_S[\widehat{\mathfrak{R}}_S(G)].$$

Therefore:
$$\mathbb{E}_S[\Psi(S)] \leq 2\mathbb{E}_S[\widehat{\mathfrak{R}}_S(G)] - 2\mathbb{E}_S[\widehat{\mathfrak{R}}_S(G)] = 0.$$
\end{proof}

With $c_i = \frac{3}{m}$ for all $i$ and $\mathbb{E}[\Psi(S)] \leq 0$:

$$\sum_{i=1}^m c_i^2 = m \cdot \frac{9}{m^2} = \frac{9}{m}.$$

By McDiarmid's inequality, with probability at least $1-\delta$:
$$\Psi(S) \leq 0 + \sqrt{\frac{9/m}{2m}\log\frac{1}{\delta}} = \sqrt{\frac{9\log(1/\delta)}{2m^2}} = \frac{3}{\sqrt{2}}\sqrt{\frac{\log(1/\delta)}{m^2}}.$$

Wait, let me recalculate. We have:
$$\Psi(S) \leq \mathbb{E}[\Psi(S)] + \sqrt{\frac{\sum_{i=1}^m c_i^2}{2}\log\frac{1}{\delta}}.$$

With $\sum_{i=1}^m c_i^2 = \frac{9}{m}$:
$$\Psi(S) \leq 0 + \sqrt{\frac{9}{2m}\log\frac{1}{\delta}} = 3\sqrt{\frac{\log(1/\delta)}{2m}}.$$

From $\Psi(S) \leq 3\sqrt{\frac{\log(1/\delta)}{2m}}$, we have:
$$\sup_{g \in G} \left[\mathbb{E}[g] - \widehat{\mathbb{E}}_S[g]\right] \leq 2\widehat{\mathfrak{R}}_S(G) + 3\sqrt{\frac{\log(1/\delta)}{2m}}.$$

Therefore, with probability at least $1-\delta$:
$${\mathbb{E}[g] \leq \widehat{\mathbb{E}}_S[g] + 2\widehat{\mathfrak{R}}_S(G) + 3\sqrt{\frac{\log(1/\delta)}{2m}}}$$

for all $g \in G$.

\textbf{Improvement:} The confidence term has $\log(1/\delta)$ instead of $\log(2/\delta)$, which is a slight improvement by a factor of $\sqrt{\log 2} \approx 0.83$ in the confidence term.

\section*{C. SVM \& Kernel Methods}

\begin{enumerate}
\item Let $\mathcal{X} \subset \mathbb{R}^n$ and $K: \mathcal{X} \times \mathcal{X} \to \mathbb{R}$ be a positive definite kernel with RKHS $\mathbb{H}$ and feature map $\phi: \mathcal{X} \to \mathbb{H}$. Consider a binary classification dataset $S = \{(x_i, y_i)\}_{i=1}^m$ with $y_i \in \{-1, 1\}$, and suppose the hard-margin SVM is feasible.

\begin{enumerate}
\item Prove that the SVM solution $w^* \in \mathbb{H}$ can be expressed as
\[
w^* = \sum_{i=1}^m \alpha_i y_i \phi(x_i), \quad \alpha_i \geq 0,
\]
and that nonzero $\alpha_i$ correspond to support vectors (\textit{Hint:} you can use the representer theorem).

\item Let $\max_i \alpha_i \leq 1$. Let $R = \sup_{x \in \mathcal{X}} \|\phi(x)\|_{\mathbb{H}}$ and $\gamma$ be the margin of the SVM. Prove that the number of support vectors $|\text{SV}|$ satisfies
\[
|\text{SV}| \geq \frac{\|w^*\|^2}{R^2} = \frac{1}{\gamma^2 R^2}.
\]

\item Let $K$ be an infinite-dimensional kernel (e.g., Gaussian RBF). Prove that there always exists a feasible hard-margin SVM for any finite dataset, and explain why the bound in part (2) is still meaningful.

\item Suppose we replace $K$ with a kernel $K_\lambda = K + \lambda I$, where $I$ is the identity in $\mathbb{H}$ and $\lambda > 0$. Prove that as $\lambda \to 0$, the solution $w_\lambda^*$ approaches the minimum-norm SVM solution, and derive an explicit bound on the number of support vectors in terms of $R$, and $\gamma$
\end{enumerate}
\end{enumerate}

\begin{enumerate}
\setcounter{enumi}{1}
\item Let $\mathcal{X} \subset \mathbb{R}^n$ and consider the function
\[
K(x, x') = \frac{1}{1 + \|x - x'\|^2} + (1 + x^\top x')^2.
\]

\begin{enumerate}
\item Prove that $K_1(x, x') = \frac{1}{1 + \|x - x'\|^2}$ is a positive definite kernel on $\mathcal{X}$. (\textit{Hint:} you can prove and use the equality $1/a = \int_0^\infty e^{-at} dt$, valid for $a > 0$.)

\item Prove that $K_2(x, x') = (1 + x^\top x')^2$ is a positive definite kernel on $\mathcal{X}$.

\item Using closure properties of PDS kernels, prove that $K = K_1 + K_2$ is positive definite.

\item Let $f: \mathbb{R} \to \mathbb{R}$ be a function with a convergent Maclaurin series $f(t) = \sum_{k=0}^\infty a_k t^k$ where $a_k \geq 0$ for all $k$. Prove that the kernel $K_f(x, x') = f(x^\top x')$ is positive definite.
\end{enumerate}
\end{enumerate}

(a) 

\textbf{Proof:}

The hard-margin SVM problem is:
\begin{align*}
\min_{w \in \mathbb{H}, b \in \mathbb{R}} \quad & \frac{1}{2}\|w\|^2 \\
\text{subject to} \quad & y_i(\langle w, \phi(x_i) \rangle + b) \geq 1, \quad \forall i \in [1, m].
\end{align*}

By the \textbf{Representer Theorem}, any solution $w^*$ to this optimization problem lies in the finite-dimensional subspace spanned by $\{\phi(x_1), \ldots, \phi(x_m)\}$. Therefore, there exist coefficients $c_1, \ldots, c_m \in \mathbb{R}$ such that
\[
w^* = \sum_{i=1}^m c_i \phi(x_i).
\]

The Lagrangian of the primal problem is:
\[
L(w, b, \alpha) = \frac{1}{2}\|w\|^2 - \sum_{i=1}^m \alpha_i [y_i(\langle w, \phi(x_i) \rangle + b) - 1],
\]
where $\alpha = (\alpha_1, \ldots, \alpha_m)$ are the Lagrange multipliers.

Since this is a convex optimization problem with linear constraints, and the problem is feasible (as given), the KKT conditions are necessary and sufficient for optimality.

The KKT conditions at the optimal point $(w^*, b^*, \alpha^*)$ are:

\begin{enumerate}
\item \textbf{Stationarity:}
\[
\nabla_w L(w^*, b^*, \alpha^*) = w^* - \sum_{i=1}^m \alpha_i^* y_i \phi(x_i) = 0
\]
\[
\nabla_b L(w^*, b^*, \alpha^*) = -\sum_{i=1}^m \alpha_i^* y_i = 0
\]

\item \textbf{Primal feasibility:}
\[
y_i(\langle w^*, \phi(x_i) \rangle + b^*) \geq 1, \quad \forall i
\]

\item \textbf{Dual feasibility:}
\[
\alpha_i^* \geq 0, \quad \forall i
\]

\item \textbf{Complementary slackness:}
\[
\alpha_i^* [y_i(\langle w^*, \phi(x_i) \rangle + b^*) - 1] = 0, \quad \forall i
\]
\end{enumerate}

From the stationarity condition, we have:
\[
w^* = \sum_{i=1}^m \alpha_i^* y_i \phi(x_i),
\]
where $\alpha_i^* \geq 0$ by the dual feasibility condition.

Defining $\alpha_i := \alpha_i^*$, we obtain the desired representation:
\[
w^* = \sum_{i=1}^m \alpha_i y_i \phi(x_i), \quad \text{with } \alpha_i \geq 0.
\]

By the complementary slackness condition, if $\alpha_i > 0$, then:
\[
y_i(\langle w^*, \phi(x_i) \rangle + b^*) - 1 = 0,
\]
which means
\[
y_i(\langle w^*, \phi(x_i) \rangle + b^*) = 1.
\]

This indicates that the point $(x_i, y_i)$ lies exactly on the margin hyperplane. By definition, such points are \textbf{support vectors}.

Conversely, if $(x_i, y_i)$ is not a support vector, then $y_i(\langle w^*, \phi(x_i) \rangle + b^*) > 1$, which by complementary slackness implies $\alpha_i = 0$.

Therefore, nonzero $\alpha_i$ correspond exactly to support vectors.

\qed

(2)

(a) 

\textbf{Proof:}  
For any $x,x'\in\mathcal{X}$, set
\[
a = 1+\|x-x'\|^2 > 0.
\]
Then
\[
K_1(x,x') = \frac{1}{1+\|x-x'\|^2} = \int_0^\infty e^{-t(1+\|x-x'\|^2)}\,dt = \int_0^\infty e^{-t} \, e^{-t\|x-x'\|^2}\,dt.
\]
For each fixed $t>0$, the function
\[
(x,x') \mapsto e^{-t\|x-x'\|^2}
\]
is a Gaussian RBF kernel and is known to be positive definite. Since a nonnegative weighted integral (with weight $e^{-t}$) of positive definite kernels is positive definite, it follows that $K_1$ is positive definite.

\qed

(b)

\textbf{Proof:}  
The kernel
\[
K_2(x,x') = \bigl(1 + x^\top x'\bigr)^2
\]
is a polynomial kernel of degree $2$. An explicit feature map can be given by
\[
\phi(x) = \Bigl(1,\sqrt{2}\, x_1, \sqrt{2}\, x_2, \dots, \sqrt{2}\, x_n, x_1^2, \sqrt{2}\, x_1 x_2, \dots, x_n^2\Bigr)^\top.
\]
Then
\[
K_2(x,x') = \langle \phi(x), \phi(x') \rangle.
\]
Since inner products in a feature space always yield positive definite kernels, $K_2$ is positive definite.

\qed

(c)

\textbf{Proof:}  
A standard closure property of positive definite (PD) kernels is that if $K^{(1)}$ and $K^{(2)}$ are PD on $\mathcal{X}$, then
\[
K(x,x') = K^{(1)}(x,x') + K^{(2)}(x,x')
\]
is also PD on $\mathcal{X}$. Since both $K_1$ and $K_2$ are PD, it immediately follows that
\[
K(x,x') = K_1(x,x') + K_2(x,x')
\]
is PD.

\qed

(d) 

\textbf{Proof:}  
For each $k \geq 0$, the kernel $(x,x') \mapsto (x^\top x')^k$ is positive definite. This can be shown:
\begin{itemize}
\item For $k=0$: constant kernel is PD.
\item For $k=1$: linear kernel $x^\top x'$ is PD.
\item For $k \geq 2$: $(x^\top x')^k$ is the product of PD kernels, hence PD by closure properties.
\end{itemize}

Since $a_k \geq 0$ and the series converges, for any finite set of points $x_1,\ldots,x_m$ and coefficients $c_1,\ldots,c_m$:
\begin{align*}
\sum_{i,j} c_i c_j K_f(x_i,x_j) &= \sum_{i,j} c_i c_j \sum_{k=0}^\infty a_k (x_i^\top x_j)^k \\
&= \sum_{k=0}^\infty a_k \sum_{i,j} c_i c_j (x_i^\top x_j)^k \geq 0,
\end{align*}
where the exchange of summation is valid by absolute convergence (for bounded $|x_i^\top x_j|$), and each inner sum is nonnegative since $(x^\top x')^k$ is PD.
\qed

\section*{D. Generalization bound based on covering numbers [Optional/Bonus problem].}

Let $H$ be a family of functions mapping $\mathcal{X}$ to a subset of real numbers $\mathcal{Y} \subseteq \mathbb{R}$. For any $\epsilon > 0$, the \textit{covering number} $\mathcal{N}(H, \epsilon)$ of $H$ for the $L_\infty$ norm is the minimal $k \in \mathbb{N}$ such that $H$ can be covered with $k$ balls of radius $\epsilon$, that is, there exists $\{h_1, \ldots, h_k\} \subseteq H$ such that, for all $h \in H$, there exists $i \leq k$ with $\|h - h_i\|_\infty = \max_{x \in \mathcal{X}} |h(x) - h_i(x)| \leq \epsilon$. In particular, when $H$ is a compact set, a finite covering can be extracted from a covering of $H$ with balls of radius $\epsilon$ and thus $\mathcal{N}(H, \epsilon)$ is finite.

Covering numbers provide a measure of the complexity of a class of functions: the larger the covering number, the richer is the family of functions. The objective of this problem is to illustrate this by proving a learning bound in the case of the squared loss. Let $D$ denote a distribution over $\mathcal{X} \times \mathcal{Y}$ according to which labeled examples are drawn. Then, the generalization error of $h \in H$ for the squared loss is defined by $R(h) = \mathbb{E}_{(x,y) \sim D}[(h(x) - y)^2]$ and its empirical error for a labeled sample $S = ((x_1, y_1), \ldots, (x_m, y_m))$ by $\widehat{R}(h) = \frac{1}{m}\sum_{i=1}^m (h(x_i) - y_i)^2$. We will assume that $H$ is bounded, that is there exists $M > 0$ such that $|h(x) - y| \leq M$ for all $(x, y) \in \mathcal{X} \times \mathcal{Y}$. The following is the generalization bound proven in this problem:
\[
\mathbb{P}_{S \sim D^m}\left[\sup_{h \in H} |R(h) - \widehat{R}(h)| \geq \epsilon\right] \leq \mathcal{N}\left(H, \frac{\epsilon}{8M}\right) 2 \exp\left(\frac{-m\epsilon^2}{2M^4}\right). \tag{1}
\]

The proof is based on the following steps.

\begin{enumerate}
\item Let $L_S = R(h) - \widehat{R}(h)$, then show that for all $h_1, h_2 \in H$ and any labeled sample $S$, the following inequality holds:
\[
|L_S(h_1) - L_S(h_2)| \leq 4M \|h_1 - h_2\|_\infty.
\]

\item Assume that $H$ can be covered by $k$ subsets $B_1, \ldots, B_k$, that is $H = B_1 \cup \ldots \cup B_k$. Then, show that, for any $\epsilon > 0$, the following upper bound holds:
\[
\mathbb{P}_{S \sim D^m}\left[\sup_{h \in H} |L_S(h)| \geq \epsilon\right] \leq \sum_{i=1}^k \mathbb{P}_{S \sim D^m}\left[\sup_{h \in B_i} |L_S(h)| \geq \epsilon\right].
\]

\item Finally, let $k = \mathcal{N}(H, \frac{\epsilon}{8M})$ and let $B_1, \ldots, B_k$ be balls of radius $\epsilon/(8M)$ centered at $h_1, \ldots, h_k$ covering $H$. Use part (a) to show that for all $i \in [1, k]$,
\[
\mathbb{P}_{S \sim D^m}\left[\sup_{h \in B_i} |L_S(h)| \geq \epsilon\right] \leq \mathbb{P}_{S \sim D^m}\left[|L_S(h_i)| \geq \frac{\epsilon}{2}\right],
\]
and apply Hoeffding's inequality to prove (3).
\end{enumerate}

(1)
We need to show that for all $h_1, h_2 \in H$ and any labeled sample $S$,
\[
|L_S(h_1) - L_S(h_2)| \leq 4M \|h_1 - h_2\|_\infty.
\]

\begin{proof}
Recall that $L_S(h) = R(h) - \widehat{R}(h)$. We have:
\begin{align*}
L_S(h_1) - L_S(h_2) &= [R(h_1) - \widehat{R}(h_1)] - [R(h_2) - \widehat{R}(h_2)]\\
&= [R(h_1) - R(h_2)] - [\widehat{R}(h_1) - \widehat{R}(h_2)].
\end{align*}

For the generalization error difference:
\begin{align*}
|R(h_1) - R(h_2)| &= \left|\mathbb{E}_{(x,y) \sim D}[(h_1(x) - y)^2 - (h_2(x) - y)^2]\right|\\
&= \left|\mathbb{E}_{(x,y) \sim D}[(h_1(x) - y + h_2(x) - y)(h_1(x) - h_2(x))]\right|\\
&= \left|\mathbb{E}_{(x,y) \sim D}[(h_1(x) + h_2(x) - 2y)(h_1(x) - h_2(x))]\right|\\
&\leq \mathbb{E}_{(x,y) \sim D}[|h_1(x) + h_2(x) - 2y| \cdot |h_1(x) - h_2(x)|]\\
&\leq \mathbb{E}_{(x,y) \sim D}[|h_1(x) - y + h_2(x) - y| \cdot |h_1(x) - h_2(x)|]\\
&\leq \mathbb{E}_{(x,y) \sim D}[(|h_1(x) - y| + |h_2(x) - y|) \cdot |h_1(x) - h_2(x)|]\\
&\leq (M + M) \|h_1 - h_2\|_\infty = 2M \|h_1 - h_2\|_\infty.
\end{align*}

Similarly, for the empirical error difference:
\begin{align*}
|\widehat{R}(h_1) - \widehat{R}(h_2)| &= \left|\frac{1}{m}\sum_{i=1}^m [(h_1(x_i) - y_i)^2 - (h_2(x_i) - y_i)^2]\right|\\
&= \left|\frac{1}{m}\sum_{i=1}^m (h_1(x_i) + h_2(x_i) - 2y_i)(h_1(x_i) - h_2(x_i))\right|\\
&\leq \frac{1}{m}\sum_{i=1}^m |h_1(x_i) + h_2(x_i) - 2y_i| \cdot |h_1(x_i) - h_2(x_i)|\\
&\leq \frac{1}{m}\sum_{i=1}^m (|h_1(x_i) - y_i| + |h_2(x_i) - y_i|) \cdot |h_1(x_i) - h_2(x_i)|\\
&\leq 2M \|h_1 - h_2\|_\infty.
\end{align*}

Therefore:
\begin{align*}
|L_S(h_1) - L_S(h_2)| &\leq |R(h_1) - R(h_2)| + |\widehat{R}(h_1) - \widehat{R}(h_2)|\\
&\leq 2M \|h_1 - h_2\|_\infty + 2M \|h_1 - h_2\|_\infty\\
&= 4M \|h_1 - h_2\|_\infty.
\end{align*}
\end{proof}

(2)

We need to show that if $H = B_1 \cup \ldots \cup B_k$, then for any $\epsilon > 0$,
\[
\mathbb{P}_{S \sim D^m}\left[\sup_{h \in H} |L_S(h)| \geq \epsilon\right] \leq \sum_{i=1}^k \mathbb{P}_{S \sim D^m}\left[\sup_{h \in B_i} |L_S(h)| \geq \epsilon\right].
\]

\begin{proof}
Note that:
\begin{align*}
\left\{\sup_{h \in H} |L_S(h)| \geq \epsilon\right\} &= \left\{\exists h \in H : |L_S(h)| \geq \epsilon\right\}\\
&= \left\{\exists h \in \bigcup_{i=1}^k B_i : |L_S(h)| \geq \epsilon\right\}\\
&= \bigcup_{i=1}^k \left\{\exists h \in B_i : |L_S(h)| \geq \epsilon\right\}\\
&= \bigcup_{i=1}^k \left\{\sup_{h \in B_i} |L_S(h)| \geq \epsilon\right\}.
\end{align*}

By the union bound (subadditivity of probability):
\[
\mathbb{P}_{S \sim D^m}\left[\bigcup_{i=1}^k \left\{\sup_{h \in B_i} |L_S(h)| \geq \epsilon\right\}\right] \leq \sum_{i=1}^k \mathbb{P}_{S \sim D^m}\left[\sup_{h \in B_i} |L_S(h)| \geq \epsilon\right].
\]
\end{proof}

(3)

Let $k = \mathcal{N}(H, \frac{\epsilon}{8M})$ and let $B_1, \ldots, B_k$ be balls of radius $\epsilon/(8M)$ centered at $h_1, \ldots, h_k$ covering $H$.

\begin{proof}
We first show that for all $i \in [1, k]$,
\[
\mathbb{P}_{S \sim D^m}\left[\sup_{h \in B_i} |L_S(h)| \geq \epsilon\right] \leq \mathbb{P}_{S \sim D^m}\left[|L_S(h_i)| \geq \frac{\epsilon}{2}\right].
\]

For any $h \in B_i$, we have $\|h - h_i\|_\infty \leq \frac{\epsilon}{8M}$. By Part 1:
\[
|L_S(h) - L_S(h_i)| \leq 4M \|h - h_i\|_\infty \leq 4M \cdot \frac{\epsilon}{8M} = \frac{\epsilon}{2}.
\]

Therefore, if $|L_S(h_i)| < \frac{\epsilon}{2}$, then for all $h \in B_i$:
\[
|L_S(h)| \leq |L_S(h_i)| + |L_S(h) - L_S(h_i)| < \frac{\epsilon}{2} + \frac{\epsilon}{2} = \epsilon.
\]

This means:
\[
\left\{|L_S(h_i)| < \frac{\epsilon}{2}\right\} \subseteq \left\{\sup_{h \in B_i} |L_S(h)| < \epsilon\right\}.
\]

Taking complements:
\[
\left\{\sup_{h \in B_i} |L_S(h)| \geq \epsilon\right\} \subseteq \left\{|L_S(h_i)| \geq \frac{\epsilon}{2}\right\}.
\]

Hence:
\[
\mathbb{P}_{S \sim D^m}\left[\sup_{h \in B_i} |L_S(h)| \geq \epsilon\right] \leq \mathbb{P}_{S \sim D^m}\left[|L_S(h_i)| \geq \frac{\epsilon}{2}\right].
\]

Apply Hoeffding's inequality to bound $\mathbb{P}_{S \sim D^m}[|L_S(h_i)| \geq \frac{\epsilon}{2}]$.

Note that $L_S(h_i) = R(h_i) - \widehat{R}(h_i) = \mathbb{E}_{(x,y)\sim D}[(h_i(x) - y)^2] - \frac{1}{m}\sum_{j=1}^m (h_i(x_j) - y_j)^2$.

Let $Z_j = (h_i(x_j) - y_j)^2$. Then $Z_j \in [0, M^2]$ (since $|h_i(x) - y| \leq M$), and:
\[
L_S(h_i) = \mathbb{E}[Z_j] - \frac{1}{m}\sum_{j=1}^m Z_j.
\]

By Hoeffding's inequality, for independent random variables $Z_1, \ldots, Z_m$ with $Z_j \in [a_j, b_j]$:
\[
\mathbb{P}\left[\left|\frac{1}{m}\sum_{j=1}^m Z_j - \mathbb{E}\left[\frac{1}{m}\sum_{j=1}^m Z_j\right]\right| \geq t\right] \leq 2\exp\left(\frac{-2m^2t^2}{\sum_{j=1}^m (b_j - a_j)^2}\right).
\]

In our case, $a_j = 0$, $b_j = M^2$, so:
\[
\mathbb{P}_{S \sim D^m}\left[|L_S(h_i)| \geq \frac{\epsilon}{2}\right] \leq 2\exp\left(\frac{-2m^2(\epsilon/2)^2}{m \cdot M^4}\right) = 2\exp\left(\frac{-m\epsilon^2}{2M^4}\right).
\]

Combine the results.

By above steps:
\begin{align*}
\mathbb{P}_{S \sim D^m}\left[\sup_{h \in H} |L_S(h)| \geq \epsilon\right] &\leq \sum_{i=1}^k \mathbb{P}_{S \sim D^m}\left[\sup_{h \in B_i} |L_S(h)| \geq \epsilon\right]\\
&\leq \sum_{i=1}^k \mathbb{P}_{S \sim D^m}\left[|L_S(h_i)| \geq \frac{\epsilon}{2}\right]\\
&\leq k \cdot 2\exp\left(\frac{-m\epsilon^2}{2M^4}\right)\\
&= \mathcal{N}\left(H, \frac{\epsilon}{8M}\right) \cdot 2\exp\left(\frac{-m\epsilon^2}{2M^4}\right).
\end{align*}

This completes the proof of inequality (1).
\end{proof}

\end{document}